% !TEX root = ../main.tex

\chapter{Implementación de la Solución}
\label{chap:implementacion-solucion}

Este capítulo describe cómo se materializó en \textsf{GHC} el diseño de la solución presentado en el capítulo anterior. Primero se introduce el flujo de compilación relevante de \textsf{GHC} y se identifican los puntos intervenidos; luego se presenta el mecanismo declarativo que permite reconocer un bloque conmutativo y se detalla su integración en el Renamer. A continuación se desarrollan la construcción del grafo de precedencia RAW y la generación recursiva de reordenamientos válidos. Finalmente, se explica cómo cada reordenamiento es procesado por \texttt{ApplicativeDo}, evaluado mediante el modelo de costo y reducido al plan que continúa en el flujo normal del compilador.

% !TEX root = ../../main.tex

\section{Flujo de compilación de GHC}
\label{sec:solucion-pipeline-ghc}

Antes de ubicar la extensión dentro de \textsf{GHC}, es necesario presentar las etapas del compilador que atraviesa una expresión de \textsf{Haskell}. La descripción se concentra en el \textit{frontend}, desde la lectura del programa hasta la producción de \texttt{Core}, porque en ese recorrido se encuentran la representación y la información utilizadas por la solución.

Al comenzar la compilación, las opciones de la línea de comandos y los pragmas \texttt{LANGUAGE} del módulo determinan las extensiones habilitadas. Esta configuración forma parte de las opciones dinámicas del compilador y es consultada por las etapas que requieren cada característica. Por ejemplo, el Parser comprueba que \texttt{QualifiedDo} permita la sintaxis de un bloque \texttt{do} calificado, mientras que el Renamer consulta \texttt{ApplicativeDo} antes de ejecutar su postprocesamiento. Por tanto, la comprobación de las extensiones no constituye una transformación separada del programa, sino una condición que guía el comportamiento de las fases posteriores.

La mayoría de los compiladores representan internamente un programa mediante un \textit{árbol de sintaxis abstracta}, o AST por su nombre en inglés. Esta estructura organiza jerárquicamente construcciones como módulos, declaraciones, expresiones y statements. Sus nodos conservan la estructura necesaria para analizar y transformar el programa, sin depender directamente de la secuencia de caracteres del archivo fuente.

En \textsf{GHC}, el AST de \textsf{Haskell} está parametrizado por la etapa de compilación. Por ejemplo, una expresión se representa mediante \texttt{HsExpr p}, donde el parámetro \texttt{p} determina la información disponible en sus nodos. Los alias \texttt{GhcPs}, \texttt{GhcRn} y \texttt{GhcTc} identifican, respectivamente, las representaciones producidas por el Parser, el Renamer y el Typechecker. No se trata de tres árboles sin relación, sino de una misma familia de sintaxis cuyos campos e invariantes evolucionan a medida que el compilador incorpora nueva información. La Expresión~\ref{expr:fases-ast-ghc} resume este recorrido para una expresión.

\begin{ReportExpression}{Fases relevantes del AST de una expresión en \textsf{GHC}.}{expr:fases-ast-ghc}
\[
\begin{aligned}
\text{código fuente}
&\xrightarrow{\mathrm{Parser}}
\texttt{HsExpr GhcPs}
\xrightarrow{\mathrm{Renamer}}
\texttt{HsExpr GhcRn}\\
&\xrightarrow{\mathrm{Typechecker}}
\texttt{HsExpr GhcTc}
\xrightarrow{\mathrm{Desugarer}}
\texttt{Core}.
\end{aligned}
\]
\end{ReportExpression}

El Parser reconoce la sintaxis del programa y construye la primera representación estructurada. En particular, un bloque \texttt{do} se incorpora como una expresión \texttt{HsDo} dentro del AST, el cual contiene el alias \texttt{GhcPs}. En esta etapa los identificadores aún corresponden a los nombres escritos en el programa y no se han resuelto sus definiciones ni sus alcances.

El Renamer transforma esa representación en un nodo del AST con alias \texttt{GhcRn}. Para ello resuelve cada identificador, determina su alcance y asigna identidades únicas de tipo \texttt{Name} a los binders. También calcula las variables libres de cada statement agrupandolas en una estructura adecuada, y ejecuta el postprocesamiento de \texttt{ApplicativeDo} sobre los bloques habilitados. Esta combinación de estructura sintáctica e información de nombres es la que permite analizar las dependencias locales sin confundir identificadores textualmente iguales que corresponden a vinculaciones distintas.

El Typechecker verifica los tipos de las expresiones y las restricciones exigidas por las clases utilizadas, y registra esa información en el mismo nodo del AST pero con alias \texttt{GhcTc}. En la solución propuesta, esta etapa comprueba indirectamente la restricción \texttt{CommutativeMonad} a través de los tipos de las operaciones calificadas empleadas por el bloque. Finalmente, el Desugarer traduce la expresión tipada a \texttt{Core}, donde las agrupaciones introducidas por \texttt{ApplicativeDo} ya se expresan mediante las operaciones aplicativas y monádicas correspondientes.

La implementación interviene en dos zonas de este recorrido. Por una parte, en la infraestructura de opciones dinámicas de \path{GHC} se incorpora la instrumentación experimental para la selección de reordenamientos válidos forzada, la cual se explicará en el capítulo de Validación. Por otra parte, la modificación funcional se concentra en el Renamer de \textsf{GHC}, donde se detecta el marcador conmutativo asociado a la declaración explicita hecha por el programador, se generan los reordenamientos válidos y se selecciona el plan de ejecución que permanecerá en el nodo del AST renombrado.

Además, se agregaron dos nuevos módulos de \textsf{Haskell} asociados a las mónadas conmutativas llamados ambos \texttt{CommutativeDo}, estos debieron agregarse en los archivos asociados a \texttt{ghc-internal} y en \texttt{base}, junto con las entradas necesarias para exponerlos desde ambas bibliotecas. El Parser, el Typechecker y el Desugarer no requirieron modificaciones: la solución reutiliza el soporte que estas etapas ya proporcionaban para \texttt{QualifiedDo} y \texttt{ApplicativeDo}. La sección siguiente desarrolla la interfaz declarativa que conecta dichos módulos con la modificación del Renamer.

% !TEX root = ../../main.tex

\section{Activación mediante CommutativeDo}
\label{sec:solucion-qualified-do}

La implementación de la solución requiere que la hipótesis de conmutatividad sea declarada de manera explícita y local a cada bloque \texttt{do}. Para proporcionar esta interfaz se incorporaron los módulos específicos \path{GHC.Internal.Control.Monad.CommutativeDo}, dentro de \texttt{ghc-internal}, y \path{Control.Monad.CommutativeDo}, dentro de \texttt{base}. El primero contiene el soporte interno, mientras que el segundo reexporta la interfaz destinada a los programas de usuario.

Dentro de estos módulos se define la clase \texttt{CommutativeMonad}, cuya única superclase es \texttt{Monad}. Una mónada utilizada mediante esta notación debe contar con una instancia de la clase. La instancia constituye una afirmación del programador acerca de la posibilidad de intercambiar efectos independientes; \textsf{GHC} no intenta demostrar la propiedad de conmutatividad. Las operaciones exportadas por el módulo incorporan la restricción \texttt{CommutativeMonad}, por lo que su cumplimiento se comprueba posteriormente durante el typechecking.

Para hacer uso de la solución dentro de un programa de \textsf{Haskell}, se utilizó la extensión \texttt{QualifiedDo}, esta permite anteponer el nombre de un módulo a un bloque \texttt{do}, por ejemplo \texttt{M.do} \cite{GHCQualifiedDo}. Este nombre actúa como calificador e indica al compilador que debe interpretar los statements utilizando las operaciones exportadas por el módulo \texttt{M}, en lugar de las operaciones monádicas habituales sin calificar. De esta manera, la interfaz puede seleccionarse localmente para un bloque, mientras que los demás bloques \texttt{do} del programa conservan su comportamiento normal.

En esta implementación, el módulo público se importa calificadamente, habitualmente como \texttt{CD}. Un bloque declarado como conmutativo se escribe comenzando con \texttt{CD.do} y construyendo su resultado terminal mediante \texttt{CD.return}. El calificador hace que los statements utilicen las operaciones exportadas por \texttt{CommutativeDo}, cuyos tipos incorporan la restricción \texttt{CommutativeMonad}. Además de las operaciones monádicas requeridas por \texttt{QualifiedDo}, el módulo exporta las operaciones aplicativas que puede necesitar la transformación posterior de \texttt{ApplicativeDo}.

Para conectar esta interfaz con el compilador, ambos módulos exportan el marcador \texttt{\_\_commutative\_do\_\_}. Durante el renombrado se consulta la presencia de este nombre en el módulo que califica \texttt{CD.do}. Esta comprobación es independiente de la resolución de la instancia: el marcador habilita la ruta de reordenamiento en el Renamer, mientras que la clase restringe las operaciones y se verifica en una etapa posterior. La activación requiere, además, que \texttt{ApplicativeDo} y \texttt{QualifiedDo} estén habilitados.

De esta forma, para aplicar la extensión experimental a un bloque \texttt{do} conmutativo como el presentado en el Código~\ref{lst:do-notation-maybe-sum}, se deben habilitar ambas extensiones, importar la interfaz calificada y declarar la instancia de \texttt{CommutativeDo} correspondiente. El Código~\ref{lst:commutative-do-maybe-sum} muestra la declaración resultante.

\noindent\begin{minipage}{\textwidth}
\begin{lstlisting}[style=haskell,caption={Uso de la extensión experimental \texttt{CommutativeDo} sobre una suma monádica con \texttt{Maybe Int}.},label={lst:commutative-do-maybe-sum}]
    {-# LANGUAGE ApplicativeDo #-}
    {-# LANGUAGE QualifiedDo #-}

    import qualified Control.Monad.CommutativeDo as CD

    instance CD.CommutativeMonad Maybe

    sumaConmutativa :: Maybe Int
    sumaConmutativa = CD.do
        x <- Just 3
        y <- Just 2
        CD.return ((+) x y)    
\end{lstlisting}
\end{minipage}

Los pragmas habilitan los dos mecanismos requeridos por la solución, mientras que la importación proporciona el calificador \texttt{CD}. La instancia declara la hipótesis de conmutatividad para \texttt{Maybe} y el uso de \texttt{CD.do} identifica localmente el bloque que seguirá la ruta experimental. La declaración conserva los cómputos del Código~\ref{lst:do-notation-maybe-sum} y produce el mismo resultado, \texttt{Just 5}.

Con estas condiciones establecidas, la sección siguiente detalla la bifurcación incorporada al flujo de renombrado.

% !TEX root = ../../main.tex

\section{Integración en el Renamer}
\label{sec:solucion-integracion-renamer}

La integración de la solución comienza en el renombramiento de toda expresión de \textsf{Haskell}, lógica implementada por la función \texttt{rnExpr}, especificamente cuando debe procesar un nodo del AST con la expresión \texttt{HsDo}. Esta debe invocar a \texttt{rnStmtsWithFreeVars} para renombrar todos los statements y asociar a cada cuerpo el conjunto de variables libres calculado durante el recorrido. A continuación, la función \texttt{postProcessStmtsForApplicativeDo} decide si debe aplicar la transformación de \texttt{ApplicativeDo} o devolver los statements sin dicho postprocesamiento.

El postprocesamiento de un bloque \texttt{do} se ejecuta cuando \texttt{ApplicativeDo} está habilitada, ingresando a la función \texttt{rearrangeForApplicativeDo}, continuando con el proceso de renombrado.

Dentro de esa función, \texttt{isCommutativeQualifiedDo} consulta si el contexto es un bloque \texttt{do} calificado y si el módulo calificador permite resolver el marcador \texttt{\_\_commutative\_do\_\_}. Cuando el marcador no se encuentra, el orden original constituye el único candidato y \texttt{ApplicativeDo} aplica su algoritmo habitual. Cuando el marcador se encuentra, la extensión genera todos los reordenamientos válidos de statements, aplica el algoritmo de \textit{rearrange} sobre cada uno y selecciona el plan de ejcución de menor costo. La Figura~\ref{fig:solucion-flujo-renamer} presenta las rutas de compilación del Renamer con la extensión integrada.

\begin{figure}[ht]
\centering
\vspace{0.2cm}
\captionsetup{skip=12pt}
\begin{tikzpicture}[
    node distance=0.48cm and 0.65cm,
    renamer step/.style={
        draw=codeframe,
        fill=codebg,
        rounded corners=2pt,
        minimum width=4.2cm,
        minimum height=0.72cm,
        text width=3.8cm,
        align=center,
        font=\fontsize{8.0pt}{10pt}\selectfont
    },
    renamer contribution/.style={
        renamer step,
        draw=codekeyword,
        fill=codekeyword!18!codebg
    },
    renamer decision/.style={
        diamond,
        draw=codekeyword,
        fill=codebg,
        aspect=2.0,
        inner sep=0.0pt,
        text width=3.2cm,
        align=center,
        font=\fontsize{8.0pt}{10pt}\selectfont
    },
    renamer arrow/.style={-latex,draw=codekeyword,line width=0.7pt},
    renamer branch/.style={fill=white,inner sep=1pt,font=\fontsize{8.0pt}{10pt}\selectfont}
]
    \node[renamer step] (hsdo) {Renombrar una expresión \texttt{HsDo} asociada a un bloque \texttt{do}};
    \node[renamer step,below=of hsdo] (rename) {Renombrar los statements y asociar sus variables libres};
    \node[renamer step,below=of rename] (postprocess) {Postprocesar los statements para \texttt{ApplicativeDo}};
    \node[renamer decision,below=0.65cm of postprocess] (ado) {¿Está habilitada \texttt{ApplicativeDo}?};
    \node[renamer step,below=0.65cm of ado] (enter) {Entrar en \texttt{rearrangeForApplicativeDo}};
    \node[renamer step,right=of enter] (noado) {Continuar sin postprocesamiento de \texttt{ApplicativeDo}};
    \node[renamer decision,below=0.65cm of enter] (marker) {¿Se resuelve el marcador conmutativo?};
    \node[renamer contribution,below=0.65cm of marker] (orders) {Generar reordenamientos válidos};
    \node[renamer step,left=of orders] (ordinary) {Aplicar \textit{rearrange} sobre el orden original};
    \node[renamer contribution,below=of orders] (rearrange) {Aplicar \textit{rearrange} a cada reordenamiento};
    \node[renamer contribution,below=of rearrange] (select) {Seleccionar el plan de menor costo};
    \node[renamer step,below=0.65cm of select] (continue) {Continuar con la etapa siguiente del compilador};

    \draw[renamer arrow] (hsdo) -- (rename);
    \draw[renamer arrow] (rename) -- (postprocess);
    \draw[renamer arrow] (postprocess) -- (ado);
    \draw[renamer arrow] (ado) -- node[renamer branch,left] {sí} (enter);
    \draw[renamer arrow] (ado.east) -| node[renamer branch,near start] {no} (noado.north);
    \draw[renamer arrow] (enter) -- (marker);
    \draw[renamer arrow] (marker) -- node[renamer branch,left] {sí} (orders);
    \draw[renamer arrow] (marker.west) -| node[renamer branch,near start] {no} (ordinary.north);
    \draw[renamer arrow] (orders) -- (rearrange);
    \draw[renamer arrow] (rearrange) -- (select);
    \draw[renamer arrow] (select) -- (continue);
    \draw[renamer arrow] (ordinary.south) |- (continue.west);
    \draw[renamer arrow] (noado.south) |- (continue.east);
\end{tikzpicture}
\caption{Flujo de renombrado y condiciones de activación de la ruta de reordenamiento de statements.}
\label{fig:solucion-flujo-renamer}
\end{figure}

Dentro de ese alcance, la Figura~\ref{fig:solucion-flujo-renamer} muestra que \texttt{ApplicativeDo} es una condición previa para ingresar en el postprocesamiento. La función \texttt{rearrangeForApplicativeDo} solo se invoca cuando esa condición se cumple; por ello, la presencia del marcador no activa por sí sola la generación de reordenamientos. Del mismo modo, \texttt{QualifiedDo} ya debió permitir la sintaxis calificada durante el parseo. Solo la combinación de un bloque calificado aceptado, \texttt{ApplicativeDo} habilitada y el marcador resoluble alcanza la ruta resaltada.

Ubicar la modificación en este punto evita realizar una reescritura externa o analizar una representación posterior que ya haya perdido la estructura del bloque \texttt{do}. También permite utilizar las identidades únicas y las variables libres producidas durante el renombrado tanto para construir las dependencias como para generar los reordenamientos válidos que recibirá \texttt{ApplicativeDo}. Las tres secciones siguientes desarrollan, en ese orden, la formación del grafo de precedencia RAW, la generación de reordenamientos válidos y la selección del plan de ejecución de menor costo.

% !TEX root = ../../main.tex

\section{Construcción del Grafo de Precedencia RAW}
\label{sec:solucion-implementacion-grafo}

La ruta resaltada en la Figura~\ref{fig:solucion-flujo-renamer} comienza aplicando el modelo de dependencias de la Sección~\ref{sec:solucion-modelo-dependencias} sobre los statements renombrados. Antes de formar el grafo, la función \texttt{rearrangeForApplicativeDo} separa el \texttt{LastStmt}, que produce el resultado del bloque \texttt{do} y permanece fijo. Considerando los demás statements $s_1,\ldots,s_n$, cada uno se asocia con su posición estable, las variables libres de su cuerpo, sus binders y las lecturas locales obtenidas al restringir esas variables libres a nombres definidos dentro del mismo bloque \texttt{do}.

En vez de introducir una estructura de datos nueva, el grafo de precedencia RAW se construye con la biblioteca \path{GHC.Data.Graph.Directed}. Cada vértice almacena la información de un statement y utiliza su posición original como clave. Para cada par $i<j$, se agrega una arista $i\to j$ cuando los nombres escritos por $s_i$ intersectan las lecturas locales de $s_j$. El Código~\ref{lst:algoritmo-construccion-grafo} muestra la construcción de este grafo utilizando pseudocódigo.

\begin{lstlisting}[style=haskell,caption={Pseudocódigo para construir el grafo de precedencia RAW.},label={lst:algoritmo-construccion-grafo}]
buildRawPrecedenceGraph(statements):
    n := length(statements)
    allWrites := union(write(s) for s in statements)
    graph := graph with vertices {1, ..., n}

    for each s_i in statements:
        writesLocal[i] := write(s_i)
        readsLocal[i]  := intersection(freeVars(s_i), allWrites)

    for i := 1 to n:
        for j := i + 1 to n:
            if intersection(writesLocal[i], readsLocal[j]) != empty:
                add edge i -> j to graph

    return graph
\end{lstlisting}

La condición $i<j$ conserva la precedencia del programa original y garantiza la aciclicidad del grafo, porque los índices aumentan estrictamente a lo largo de cada arista. La construcción de este realiza

\[
\sum_{i=1}^{n-1}(n-i)=\frac{n(n-1)}{2}
\]

comparaciones entre pares. Por tanto, la cantidad de pruebas de dependencia es $O(n^2)$, sin considerar el costo interno de construir e intersectar los conjuntos de nombres. En el ejemplo probabilístico del Código~\ref{lst:main-example}, la construcción del grafo de precedencia RAW produce las aristas $1\to2$, $1\to4$ y $3\to4$, presentadas en la Figura~\ref{fig:main-example-precedence-graph}.

% !TEX root = ../../main.tex

\section{Generación de Reordenamientos Válidos}
\label{sec:solucion-implementacion-permutaciones}

La función \texttt{generateAllSemanticTopSorts} recibe el grafo de precedencia RAW y materializa los reordenamientos válidos como sus ordenamientos topológicos. Para ello calcula el grado de entrada de cada vértice mediante el grafo original y su transposición. Un vértice cuyo grado de entrada es cero no posee predecesores, por lo que puede ser declarado para ser el siguiente statement en el reordenamiento calculado. Luego de ser declarado, todas las aristas salientes de este mismo son eliminadas para así garantizar que los statements originalmente dependendientes de este puedan ser declarados y así continuar con la generación del ordenamiento topológico.

El algoritmo de Kahn clásico selecciona uno a uno los vértices del grafo de precedencia y produce un único ordenamiento topológico \cite{Kahn1962}. La implementación de la solución reemplaza esa elección por una ramificación sobre todos los vértices disponibles. Cada rama emite un vértice, decrementa el grado de entrada de sus sucesores y repite el procedimiento sobre el estado resultante. Cuando la secuencia acumulada de statements contiene todos los vértices del grafo, se obtiene un orden topológico y por ende un reordenamiento válido. El Código~\ref{lst:algoritmo-enumeracion-topologica} presenta el procedimiento abstracto.

\newpage

\begin{lstlisting}[style=haskell,caption={Pseudocódigo del algoritmo que genera todos los ordenamientos topológicos del grafo de precendencia RAW de un bloque \texttt{do}.},label={lst:algoritmo-enumeracion-topologica}]
generateAllSemanticTopSorts(graph):
    indegree := computeIndegrees(graph)
    extendTopSorts(graph, indegree, empty, [])

extendTopSorts(graph, indegree, emitted, order):
    if length(order) == numberOfVertices(graph):
        emit order
        return

    choices := every non-emitted vertex v with indegree[v] == 0

    for each v in choices:
        nextIndegree := copy(indegree)
        for each successor w of v:
            nextIndegree[w] := nextIndegree[w] - 1

        extendTopSorts(
            graph,
            nextIndegree,
            union(emitted, {v}),
            order ++ [v])
\end{lstlisting}

Cada llamada recursiva agrega únicamente un vértice cuyos predecesores ya fueron emitidos, por lo que todo resultado completo conserva las aristas del grafo. A la inversa, la primera posición de cualquier ordenamiento topológico debe tener grado de entrada cero; como el algoritmo ramifica sobre todas esas posibilidades y repite el criterio sobre el grafo restante, también alcanza cada reordenamiento válido.

Los índices se recorren en orden creciente para producir resultados estables. En particular, la primera ramificación completa del algoritmo reproduce el orden original, que queda asociado al índice cero. Si $P$ es la cantidad de ordenamientos topológicos, la función materializa $P$ secuencias de longitud $n$. En el caso extremo de un grafo sin aristas, todos los statements pueden ocupar cualquier posición y se generan $n!$ reordenamientos válidos.

La Figura~\ref{fig:arbol-enumeracion-main-example} muestra la ejecución recursiva del algoritmo sobre el bloque \texttt{do} del ejemplo probabilístico con Código~\ref{lst:main-example}. Cada estado contiene el prefijo emitido $\pi$ y el subgrafo $G[R]$ inducido por los vértices restantes. Una arista representa la elección del statement emitido y las hojas corresponden a los reordenamientos completos.

\begin{figure}[ht]
\centering
\begin{tikzpicture}[
    enumeration state/.style={
        draw=graphNodeBorder,
        fill=graphNodeBg,
        text=graphNodeFg,
        rounded corners=2pt,
        line width=0.6pt,
        minimum width=2.45cm,
        minimum height=0.9cm,
        align=center,
        font=\scriptsize
    },
    enumeration leaf/.style={
        enumeration state,
        fill=codekeyword!16!codebg
    },
    enumeration edge/.style={
        -latex,
        draw=graphEdge,
        line width=0.7pt
    },
    branch label/.style={
        fill=white,
        inner sep=1pt,
        font=\scriptsize
    }
]
    \node[enumeration state] (root) at (0,8) {$\pi=[]$\\$G[\{1,2,3,4\}]$};

    \node[enumeration state] (p1) at (-3.75,6) {$\pi=[1]$\\$G[\{2,3,4\}]$};
    \node[enumeration state] (p3) at (4.5,6) {$\pi=[3]$\\$G[\{1,2,4\}]$};

    \node[enumeration state] (p12) at (-6,4) {$\pi=[1,2]$\\$G[\{3,4\}]$};
    \node[enumeration state] (p13) at (-1.5,4) {$\pi=[1,3]$\\$G[\{2,4\}]$};
    \node[enumeration state] (p31) at (4.5,4) {$\pi=[3,1]$\\$G[\{2,4\}]$};

    \node[enumeration state] (p123) at (-6,2) {$\pi=[1,2,3]$\\$G[\{4\}]$};
    \node[enumeration state] (p132) at (-3,2) {$\pi=[1,3,2]$\\$G[\{4\}]$};
    \node[enumeration state] (p134) at (0,2) {$\pi=[1,3,4]$\\$G[\{2\}]$};
    \node[enumeration state] (p312) at (3,2) {$\pi=[3,1,2]$\\$G[\{4\}]$};
    \node[enumeration state] (p314) at (6,2) {$\pi=[3,1,4]$\\$G[\{2\}]$};

    \node[enumeration leaf] (c0) at (-6,0) {$R_0$\\$\pi=[1,2,3,4]$\\$G[\varnothing]$};
    \node[enumeration leaf] (c1) at (-3,0) {$R_1$\\$\pi=[1,3,2,4]$\\$G[\varnothing]$};
    \node[enumeration leaf] (c2) at (0,0) {$R_2$\\$\pi=[1,3,4,2]$\\$G[\varnothing]$};
    \node[enumeration leaf] (c3) at (3,0) {$R_3$\\$\pi=[3,1,2,4]$\\$G[\varnothing]$};
    \node[enumeration leaf] (c4) at (6,0) {$R_4$\\$\pi=[3,1,4,2]$\\$G[\varnothing]$};

    \draw[enumeration edge] (root) -- node[branch label] {$s_1$} (p1);
    \draw[enumeration edge] (root) -- node[branch label] {$s_3$} (p3);

    \draw[enumeration edge] (p1) -- node[branch label] {$s_2$} (p12);
    \draw[enumeration edge] (p1) -- node[branch label] {$s_3$} (p13);
    \draw[enumeration edge] (p3) -- node[branch label] {$s_1$} (p31);

    \draw[enumeration edge] (p12) -- node[branch label] {$s_3$} (p123);
    \draw[enumeration edge] (p13) -- node[branch label] {$s_2$} (p132);
    \draw[enumeration edge] (p13) -- node[branch label] {$s_4$} (p134);
    \draw[enumeration edge] (p31) -- node[branch label] {$s_2$} (p312);
    \draw[enumeration edge] (p31) -- node[branch label] {$s_4$} (p314);

    \draw[enumeration edge] (p123) -- node[branch label] {$s_4$} (c0);
    \draw[enumeration edge] (p132) -- node[branch label] {$s_4$} (c1);
    \draw[enumeration edge] (p134) -- node[branch label] {$s_2$} (c2);
    \draw[enumeration edge] (p312) -- node[branch label] {$s_4$} (c3);
    \draw[enumeration edge] (p314) -- node[branch label] {$s_2$} (c4);
\end{tikzpicture}
\caption{Árbol de llamadas para generar todos los reordenamientos válidos del programa probabilístico con Código~\ref{lst:main-example}. La raíz representa la invocación inicial y los otros 15 estados corresponden a llamadas recursivas.}
\label{fig:arbol-enumeracion-main-example}
\end{figure}

Las cinco hojas de la figura corresponden, en el mismo orden, a los ordenamientos topológicos presentados en la Tabla~\ref{tab:solucion-candidatos-ejemplo}. La hoja $R_0$ conserva la secuencia original de statements; las demás intercalan statements independientes sin invertir ninguna de las tres aristas del grafo de precedencia RAW de la Figura~\ref{fig:main-example-precedence-graph}.

% !TEX root = ../../main.tex

\section{Generación y Selección de Planes}
\label{sec:solucion-integracion-applicative-do}

Después de la generación de todos los reordenamientos válidos de un bloque \texttt{do}, como muestra la Figura~\ref{fig:solucion-flujo-renamer} el siguiente paso de la solución consiste en aplicar el algoritmo de \textit{rearrange} sobre cada uno de ellos. Por ende, cada ordenamiento topológico del grafo de precedencia RAW se reconstruye como una secuencia de statements renombrados con sus anotaciones de variables libres y se entrega a la lógica existente de \texttt{ApplicativeDo}.

Todo plan de ejecución resultante se representa programáticamente mediante el tipo \texttt{StmtTree}. Sus tres constructores, mostrados en el Código~\ref{lst:stmt-tree}, expresan un statement elemental, una composición secuencial mediante \textit{bind} o una agrupación aplicativa.

\newpage

\begin{lstlisting}[style=haskell,caption={Representación de los planes de ejecución generados por \textsf{ApplicativeDo}.},label={lst:stmt-tree}]
data StmtTree a
  = StmtTreeOne a
  | StmtTreeBind (StmtTree a) (StmtTree a)
  | StmtTreeApplicative [StmtTree a]
\end{lstlisting}

La extensión obtiene un árbol por cada reordenamiento válido. Si $P$ es la cantidad de secuencias generadas y cada una contiene $n$ statements, la etapa de \textit{rearrange} se ejecuta $P$ veces utilizando la estrategia configurada en \textsf{GHC}. Por tanto, según la complejidad establecida en la Expresión~\ref{expr:complejidad-rearrange}, la generación de todos los planes de ejecución tendrá complejidad $P*\,T_{\mathrm{rearrange}}(n)$. Esto corresponde a $O(P*n^2)$ cuando se utiliza la variante heurística predeterminada y a $O(P*n^3)$ cuando se activa la variante óptima.

Para comparar los planes se incorporó la función \texttt{stmtTreeCost}, presentada en el Código~\ref{lst:stmt-tree-cost}. La función implementa la Expresión~\ref{expr:modelo-costos-estructural}: asigna costo unitario a un statement, suma las composiciones secuenciales y adopta el máximo de las ramas aplicativas.

\begin{lstlisting}[style=haskell,caption={Cálculo del costo estructural de un plan de ejecución.},label={lst:stmt-tree-cost}]
stmtTreeCost :: StmtTree a -> Cost
stmtTreeCost (StmtTreeOne _) = 1
stmtTreeCost (StmtTreeBind l r) =
  stmtTreeCost l + stmtTreeCost r
stmtTreeCost (StmtTreeApplicative ts) =
  case ts of
    [] -> 0
    _  -> Partial.maximum (map stmtTreeCost ts)
\end{lstlisting}

La selección del mejor plan de ejecución utiliza \texttt{minimumBy} sobre la lista de arboles aplicando la función \texttt{stmtTreeCost}, escoginedo el primer mínimo, por lo que el orden estable de los reordenemientos válidos proporciona un desempate determinista. En el ejemplo probabilístico con Código~\ref{lst:main-example}, el orden original genera un arbol con costo 3 y los demás reordenamientos válidos obtienen arboles con costo 2; luego se selecciona el primero de ellos, asociado al plan de ejecución $(s_1\mid s_3);(s_2\mid s_4)$ como se muestra en la Tabla~\ref{tab:main-example-plans}.

Finalmente, la función \texttt{stmtTreeToStmts} reconstruye la lista, genera nodos \texttt{ApplicativeStmt} y reincorpora el statement terminal. Desde ese punto la extensión deja de intervenir y, como resume la Expresión~\ref{expr:fases-ast-ghc}, el nodo AST renombrado continúa por el Typechecker y el Desugarer de forma normal.

